\documentclass[10pt]{letter}
\usepackage[T1]{fontenc}
\usepackage{lmodern}
\usepackage{amsmath}
\usepackage{microtype}
\usepackage[margin=0.85in]{geometry}
\usepackage[hidelinks]{hyperref}

\signature{Siwei Zeng\\Independent Researcher\\Beijing, China\\
\href{mailto:endlesslethe@gmail.com}{endlesslethe@gmail.com}\\
ORCID: \href{https://orcid.org/0009-0008-8914-7143}{0009-0008-8914-7143}}
\address{Siwei Zeng\\Independent Researcher\\Beijing, China\\
\href{mailto:endlesslethe@gmail.com}{endlesslethe@gmail.com}}
\date{21 August 2026}

\begin{document}

\begin{letter}{Professor Vanni Noferini\\
Editor-in-Chief\\
\textit{The Electronic Journal of Linear Algebra}}

\opening{Dear Professor Noferini,}

Please consider my manuscript, ``The General Subgroup
Permanental-Dominance Conjecture in Order Four,'' for publication in
\textit{The Electronic Journal of Linear Algebra}.

The manuscript proves the complete order-four case of the general subgroup
permanental-dominance conjecture: for every subgroup $H\leq S_4$, every
irreducible complex character $\chi$ of $H$, and every $4\times4$ Hermitian
positive-semidefinite matrix $A$, it establishes
$d_\chi^H(A)/\chi(1)\leq \operatorname{per} A$.  This goes beyond the usual
$S_4$ immanant specialization by treating all thirty-seven
irreducible-character cases arising from the eleven conjugacy classes of
subgroups of $S_4$.

The main conceptual contribution is the treatment of the two non-real
$A_4$ characters.  Their inequalities are reduced to polynomial
nonnegativity on the cone of $3\times3$ positive-semidefinite Gram matrices.
A rank-one sum-of-squares identity, an exact interior certificate, and
closure prove the required positivity.  The remaining thirty-five cases are
resolved by general matrix inequalities.  A companion Lean formalization
checks the central identities and the complete finite reduction.

The general subgroup conjecture was previously known only through matrix
order three.  The present result closes the next order completely, and the
Gram-cone reduction supplies a geometric mechanism that may be useful for
related matrix inequalities.  The result and its methods lie directly within
the matrix-analysis and linear-algebra scope of the journal.

The manuscript is original, has not been published previously, and is not
under consideration by another journal.  The cited repository contains the
preprint and companion formalization.  No external funding was received, and
I have no competing interests to declare.

Thank you for considering the manuscript.

\closing{Sincerely,}

\end{letter}
\end{document}
